% =========================================================
%  ppt/01-pengantar-python-pygame.tex
%  Materi 1 — Pengantar Pemrograman Game dan Persiapan Python/Pygame
%  Bab III buku ajar
% =========================================================
\documentclass[aspectratio=169,10pt]{beamer}
% =========================================================
%  ppt/config/packages.tex
%  Paket bersama untuk seluruh presentasi kuliah
%  Pengantar Pemrograman Game
% =========================================================

\usepackage[T1]{fontenc}
\usepackage[utf8]{inputenc}
\usepackage[indonesian]{babel}
\usepackage{lmodern}
\usepackage{textcomp}
\usepackage{amsmath}
\usepackage{amssymb}
\usepackage{graphicx}
\usepackage{booktabs}
\usepackage{tabularx}
\usepackage{array}
\usepackage{multicol}
\usepackage{listings}
\usepackage{upquote}
\usepackage{adjustbox}
\usepackage{tcolorbox}
\usepackage{tikz}
\usetikzlibrary{
  arrows.meta,
  positioning,
  shapes.geometric,
  shapes.misc,
  calc,
  fit,
  backgrounds,
  matrix,
  chains,
  decorations.pathreplacing,
  patterns
}

\tcbuselibrary{skins}

% Jangan memuat geometry: ukuran halaman diatur oleh beamer.

% =========================================================
%  ppt/config/theme.tex
%  Satu tema Beamer untuk keempat belas presentasi
%  Pengantar Pemrograman Game
% =========================================================

% --------------------------------------------------------- warna
\definecolor{ubNavy}{HTML}{1F3864}
\definecolor{ubAccent}{HTML}{2E75B6}
\definecolor{ubHead}{HTML}{D9E2F3}
\definecolor{ubSoft}{HTML}{F2F6FC}
\definecolor{ubLine}{HTML}{BFBFBF}
\definecolor{ubText}{HTML}{1F2937}
\definecolor{ubMuted}{HTML}{5A6678}
\definecolor{ubGreen}{HTML}{1E7B45}
\definecolor{ubGreenBg}{HTML}{F1F8F2}
\definecolor{ubRed}{HTML}{B3261E}
\definecolor{ubRedBg}{HTML}{FDF2F1}
\definecolor{ubWarn}{HTML}{A35A05}
\definecolor{ubWarnBg}{HTML}{FFF8EC}
\definecolor{ubGray}{HTML}{F6F6F6}
\definecolor{ubCodeBg}{HTML}{10233F}
\definecolor{ubCodeText}{HTML}{E8EEF7}
\definecolor{ubCodeKey}{HTML}{8FBCF0}
\definecolor{ubCodeCom}{HTML}{8FB37A}
\definecolor{ubCodeStr}{HTML}{D9A46B}
\definecolor{ubCodeNum}{HTML}{7E99BE}

% Nama warna yang dipakai gambar TikZ buku ajar.
\definecolor{ubPrimary}{HTML}{1F3864}
\definecolor{ubAlt}{HTML}{F2F6FC}

% --------------------------------------------------------- dasar
\usetheme{default}
\usefonttheme{professionalfonts}
\setbeamertemplate{navigation symbols}{}
\setbeamerfont{frametitle}{size=\Large,series=\bfseries}
\setbeamerfont{framesubtitle}{size=\small}
\setbeamerfont{footline}{size=\tiny}
\setbeamerfont{block title}{size=\small,series=\bfseries}
\setbeamerfont{block body}{size=\small}
\setbeamerfont{itemize/enumerate body}{size=\small}
\setbeamerfont{itemize/enumerate subbody}{size=\footnotesize}

\setbeamersize{text margin left=0.60cm,text margin right=0.60cm}
\uchyph=0

\setbeamercolor{normal text}{fg=ubText,bg=white}
\setbeamercolor{structure}{fg=ubAccent}
\setbeamercolor{frametitle}{fg=ubNavy,bg=white}
\setbeamercolor{title}{fg=ubNavy}
\setbeamercolor{footline}{fg=ubMuted,bg=white}
\setbeamercolor{itemize item}{fg=ubAccent}
\setbeamercolor{itemize subitem}{fg=ubAccent}
\setbeamercolor{block title}{fg=white,bg=ubNavy}
\setbeamercolor{block body}{fg=ubText,bg=ubSoft}
\setbeamercolor{block title alerted}{fg=white,bg=ubWarn}
\setbeamercolor{block body alerted}{fg=ubText,bg=ubWarnBg}
\setbeamercolor{block title example}{fg=white,bg=ubGreen}
\setbeamercolor{block body example}{fg=ubText,bg=ubGreenBg}
\setbeamercolor{caption name}{fg=ubMuted}

\setbeamertemplate{blocks}[rounded][shadow=false]
\setbeamertemplate{itemize item}{\raisebox{0.10em}{\color{ubAccent}\scriptsize$\blacksquare$}}
\setbeamertemplate{itemize subitem}{\raisebox{0.06em}{\color{ubAccent}\tiny$\blacktriangleright$}}
\setbeamertemplate{enumerate item}{%
  \makebox[1.3em][l]{\color{ubAccent}\insertenumlabel.}}

\addtobeamertemplate{itemize/enumerate body begin}{%
  \setlength{\itemsep}{0.42em}\setlength{\parskip}{0pt}}{}

% --------------------------------------------------------- judul frame
\setbeamertemplate{frametitle}[default][left]

% --------------------------------------------------------- kaki
\newcommand{\materino}{00}
\newcommand{\materifoot}{Pengantar Pemrograman Game \textperiodcentered\ Materi \materino}

\setbeamertemplate{footline}{%
  \begin{beamercolorbox}[wd=\paperwidth,ht=0.015cm,dp=0cm]{footline}%
    \color{ubHead}\rule{\paperwidth}{0.4pt}%
  \end{beamercolorbox}%
  \begin{beamercolorbox}[wd=\paperwidth,ht=0.40cm,dp=0.14cm,%
                         leftskip=0.60cm,rightskip=0.60cm]{footline}%
    \usebeamerfont{footline}\color{ubMuted}\materifoot\hfill\insertframenumber%
  \end{beamercolorbox}%
}

% --------------------------------------------------------- makro buku ajar
\providecommand{\code}[1]{\texttt{\detokenize{#1}}}
\providecommand{\istilah}[1]{\emph{#1}}
\providecommand{\CPMK}[1]{CPMK\,#1}
\providecommand{\CPL}[1]{CPL\,#1}
\providecommand{\SUBCPMK}[1]{Sub-CPMK\,#1}

% --------------------------------------------------------- gaya kode
\lstdefinelanguage{CSharp}{
  morekeywords={abstract,as,base,bool,break,byte,case,catch,char,class,const,
    continue,decimal,default,delegate,do,double,else,enum,event,explicit,
    extern,false,finally,fixed,float,for,foreach,goto,if,implicit,in,int,
    interface,internal,is,lock,long,namespace,new,null,object,operator,out,
    override,params,private,protected,public,readonly,ref,return,sbyte,
    sealed,short,sizeof,stackalloc,static,string,struct,switch,this,throw,
    true,try,typeof,uint,ulong,unchecked,unsafe,ushort,using,virtual,void,
    volatile,while,var,get,set},
  sensitive=true,
  morecomment=[l]{//},
  morecomment=[s]{/*}{*/},
  morestring=[b]",
}

\lstdefinestyle{kode}{
  basicstyle=\ttfamily\footnotesize\color{ubCodeText},
  keywordstyle=\color{ubCodeKey}\bfseries,
  commentstyle=\color{ubCodeCom}\itshape,
  stringstyle=\color{ubCodeStr},
  numberstyle=\tiny\color{ubCodeNum},
  numbers=left,
  numbersep=7pt,
  backgroundcolor=\color{ubCodeBg},
  frame=none,
  showstringspaces=false,
  breaklines=false,
  tabsize=4,
  columns=fullflexible,
  keepspaces=true,
  xleftmargin=12pt,
  xrightmargin=2pt,
  aboveskip=0pt,
  belowskip=0pt,
  upquote=true,
}
\lstset{style=kode}

\newcommand{\kodejelas}[1]{%
  \vspace{0.16cm}%
  {\footnotesize\color{ubMuted}#1\par}%
}
\newcommand{\sumber}[1]{%
  \vspace{0.08cm}%
  {\scriptsize\color{ubMuted}\itshape #1\par}%
}
\newcommand{\awalan}[1]{%
  {\small\color{ubMuted}#1\par}\vspace{0.18cm}%
}

% --------------------------------------------------------- kotak
\newtcolorbox{kotak}[1][]{enhanced, colback=ubSoft, colframe=ubAccent,
  boxrule=0pt, leftrule=2.2pt, arc=1pt, left=6pt, right=6pt, top=4pt,
  bottom=4pt, fontupper=\footnotesize, #1}
\newtcolorbox{peringatan}[1][]{enhanced, colback=ubWarnBg, colframe=ubWarn,
  boxrule=0pt, leftrule=2.2pt, arc=1pt, left=6pt, right=6pt, top=4pt,
  bottom=4pt, fontupper=\footnotesize, #1}
\newtcolorbox{kiat}[1][]{enhanced, colback=ubGreenBg, colframe=ubGreen,
  boxrule=0pt, leftrule=2.2pt, arc=1pt, left=6pt, right=6pt, top=4pt,
  bottom=4pt, fontupper=\footnotesize, #1}

% --------------------------------------------------------- sampul
% \deckcover{nomor}{minggu}{judul}{sub-cpmk}{teknologi}{sumber}{blok}
\newcommand{\deckcover}[7]{%
\begin{frame}[plain]
\begin{tikzpicture}[remember picture, overlay]
  \fill[ubNavy] (current page.south west) rectangle ([xshift=6.00cm]current page.north west);
  \fill[ubAccent] ([xshift=6.00cm]current page.south west) rectangle ([xshift=6.08cm]current page.north west);

  \node[anchor=north west, text width=4.72cm, align=left, inner sep=0pt]
    at ($(current page.north west)+(0.66cm,-0.60cm)$) {%
      {\color{white}\bfseries\fontsize{10.2}{12.2}\selectfont
       Universitas Bale Bandung}\\[4pt]
      {\color{ubHead}\fontsize{7.6}{9.4}\selectfont
       Fakultas Teknologi Informasi}\\[1pt]
      {\color{ubHead}\fontsize{7.6}{9.4}\selectfont
       Program Studi Teknik Informatika}\\[7pt]
      {\color{ubAccent}\rule{4.0cm}{0.5pt}}%
    };

  \node[anchor=north west, text width=4.72cm, align=left, inner sep=0pt]
    at ($(current page.north west)+(0.66cm,-2.72cm)$) {%
      {\color{ubAccent}\fontsize{6.4}{7.8}\selectfont MATA KULIAH}\\[3pt]
      {\color{white}\bfseries\fontsize{9.8}{12.0}\selectfont
       Pengantar Pemrograman Game}\\[3pt]
      {\color{ubHead}\fontsize{7.6}{9.4}\selectfont
       MWCA009 \textperiodcentered\ 3 SKS}\\[8pt]
      {\color{ubHead}\fontsize{7.6}{9.4}\selectfont
       Semester Ganjil 2026/2027}%
    };

  \node[anchor=south west, text width=4.72cm, align=left, inner sep=0pt]
    at ($(current page.south west)+(0.66cm,0.52cm)$) {%
      {\color{ubAccent}\fontsize{6.4}{7.8}\selectfont DOSEN PENGAMPU}\\[3pt]
      {\color{white}\fontsize{8.4}{10.4}\selectfont
       Cecep Suwanda, S.Si., M.Kom.}%
    };

  \node[anchor=north west, text width=8.95cm, align=left, inner sep=0pt]
    at ($(current page.north east)+(-9.40cm,-0.92cm)$) {%
      {\color{ubAccent}\bfseries\fontsize{8.2}{10.2}\selectfont
       MATERI #1\ \ \textperiodcentered\ \ MINGGU #2}\\[8pt]
      {\color{ubNavy}\bfseries\fontsize{19.4}{23.4}\selectfont
       \hyphenpenalty=10000 \exhyphenpenalty=10000 \raggedright #3}%
    };

  \node[anchor=north west, text width=8.95cm, align=left, inner sep=0pt]
    at ($(current page.north east)+(-9.40cm,-4.30cm)$) {%
      \renewcommand{\arraystretch}{1.30}%
      \begin{tabular}{@{}>{\color{ubMuted}\fontsize{7.3}{9.0}\selectfont}p{1.85cm}
                        >{\color{ubText}\fontsize{8.1}{10.0}\selectfont}p{6.7cm}@{}}
        Sub-CPMK      & #4 \\
        Teknologi     & #5 \\
        Sumber materi & #6 \\
        Blok          & #7 \\
      \end{tabular}%
    };
\end{tikzpicture}
\end{frame}%
}

% --------------------------------------------------------- pembatas bagian
\newcommand{\sectionframe}[3]{%
\begin{frame}[plain]
\begin{tikzpicture}[remember picture, overlay]
  \fill[ubNavy] (current page.south west) rectangle (current page.north east);
  \fill[ubAccent] (current page.north west) rectangle ([xshift=0.10cm]current page.south west);
  \node[anchor=west, inner sep=0pt]
    at ($(current page.west)+(1.00cm,1.45cm)$)
    {\color{ubNavy!52!white}\bfseries\fontsize{58}{58}\selectfont #1};
  \node[anchor=west, inner sep=0pt, text width=11.6cm, align=left]
    at ($(current page.west)+(1.05cm,-0.55cm)$) {%
      {\color{ubAccent}\bfseries\fontsize{8.0}{10.0}\selectfont BAGIAN}\\[3pt]
      {\color{white}\bfseries\fontsize{20.5}{24.5}\selectfont #2}\\[5pt]
      {\color{ubHead}\small #3}%
    };
  \node[anchor=south east, inner sep=0pt]
    at ($(current page.south east)+(-0.60cm,0.32cm)$)
    {\color{ubHead}\fontsize{7.4}{9}\selectfont #1 / 6
     \ \ \textperiodcentered\ \ \insertframenumber};
\end{tikzpicture}
\end{frame}%
}

% --------------------------------------------------------- peta pertemuan
\newcommand{\blokpeta}[6]{%
  \ifnum#2=#3
    \node[rounded corners=3pt, fill=ubNavy, draw=ubNavy, line width=0.5pt,
          text width=3.02cm, minimum height=1.62cm, align=left, inner sep=5pt]
      at (#1,0) {%
        {\color{ubAccent}\bfseries\fontsize{6.8}{8}\selectfont #4}\\[1pt]
        {\color{white}\bfseries\fontsize{7.6}{9}\selectfont #5}\\[2pt]
        {\color{ubHead}\fontsize{6.4}{7.6}\selectfont #6}};
  \else
    \node[rounded corners=3pt, fill=ubSoft, draw=ubHead, line width=0.5pt,
          text width=3.02cm, minimum height=1.62cm, align=left, inner sep=5pt]
      at (#1,0) {%
        {\color{ubAccent}\bfseries\fontsize{6.8}{8}\selectfont #4}\\[1pt]
        {\color{ubNavy}\bfseries\fontsize{7.6}{9}\selectfont #5}\\[2pt]
        {\color{ubMuted}\fontsize{6.4}{7.6}\selectfont #6}};
  \fi
}

% \petapertemuan{blok aktif}{sebelumnya}{sesudahnya}{keluaran}
% \petapertemuan{blok aktif}{sebelumnya}{pertemuan ini}{sesudahnya}{keluaran}
\newcommand{\petapertemuan}[5]{%
\begin{frame}{Peta Pertemuan}
{\scriptsize\color{ubMuted}Kedudukan materi dalam empat blok perkuliahan}\par
\vspace{0.10cm}
\centering
\begin{tikzpicture}
  \blokpeta{-5.58}{1}{#1}{BLOK I}{Materi 1--5}{Game loop, vektor,\\fisika, collision}
  \blokpeta{-1.86}{2}{#1}{BLOK II}{Materi 6--9}{AI NPC: chase, FSM,\\pathfinding A*}
  \blokpeta{1.86}{3}{#1}{BLOK III}{Materi 10--11}{Unity: scene,\\level, physics}
  \blokpeta{5.58}{4}{#1}{BLOK IV}{Materi 12--14}{Aset visual, audio,\\UI, playtest}
\end{tikzpicture}
\par\vspace{0.30cm}
\begin{columns}[T,onlytextwidth]
  \begin{column}{0.315\textwidth}
    \begin{beamercolorbox}[wd=\linewidth,sep=4pt,rounded=true]{block body}
      {\color{ubMuted}\fontsize{6.6}{8}\selectfont SEBELUMNYA}\\[2pt]
      {\footnotesize #2}
    \end{beamercolorbox}
  \end{column}
  \begin{column}{0.315\textwidth}
    \begin{beamercolorbox}[wd=\linewidth,sep=4pt,rounded=true]{block title}
      {\color{ubHead}\fontsize{6.6}{8}\selectfont PERTEMUAN INI}\\[2pt]
      {\footnotesize #3}
    \end{beamercolorbox}
  \end{column}
  \begin{column}{0.315\textwidth}
    \begin{beamercolorbox}[wd=\linewidth,sep=4pt,rounded=true]{block body}
      {\color{ubMuted}\fontsize{6.6}{8}\selectfont SESUDAHNYA}\\[2pt]
      {\footnotesize #4}
    \end{beamercolorbox}
  \end{column}
\end{columns}
\par\vspace{0.26cm}
{\footnotesize\color{ubText}\textbf{Keluaran pertemuan.} #5\par}
\end{frame}%
}

% --------------------------------------------------------- penutup
\newcommand{\penutup}[2]{%
\begin{frame}[plain]
\begin{tikzpicture}[remember picture, overlay]
  \fill[ubNavy] (current page.south west) rectangle (current page.north east);
  \fill[ubAccent] (current page.north west) rectangle ([xshift=0.10cm]current page.south west);

  \node[anchor=north west, text width=13.6cm, align=left, inner sep=0pt]
    at ($(current page.north west)+(1.00cm,-0.95cm)$) {%
      {\color{ubAccent}\bfseries\fontsize{8.2}{10}\selectfont ASESMEN MINGGU INI}\\[6pt]
      {\color{white}\bfseries\fontsize{24}{28}\selectfont #1}%
    };

  \node[anchor=north west, text width=13.6cm, align=left, inner sep=0pt]
    at ($(current page.north west)+(1.00cm,-4.55cm)$) {%
      {\color{ubNavy!0!ubHead}\rule{13.4cm}{0.4pt}}\\[7pt]
      {\color{ubAccent}\bfseries\fontsize{7.4}{9}\selectfont
       PENGAIT KE MATERI BERIKUTNYA}\\[4pt]
      {\color{ubHead}\fontsize{10.6}{13.2}\selectfont #2}%
    };

  \node[anchor=south east, inner sep=0pt]
    at ($(current page.south east)+(-0.60cm,0.32cm)$)
    {\color{ubHead}\fontsize{7.4}{9}\selectfont \insertframenumber};
\end{tikzpicture}
\end{frame}%
}



\renewcommand{\materino}{01}

\begin{document}

\deckcover{01}{1}
  {Pengantar Pemrograman Game dan Persiapan Python/Pygame}
  {menuju 1.1}
  {Python / Pygame}
  {Bab III buku ajar}
  {Blok I \textperiodcentered\ Materi 1--5}

% --------------------------------------------------------- capaian
\begin{frame}{Capaian Pertemuan}
\setbeamerfont{itemize/enumerate body}{size=\scriptsize}
\awalan{Sub-CPMK dan indikator dari RPS, serta tujuan Bab III buku ajar.}
\begin{columns}[T,onlytextwidth]
  \begin{column}{0.52\textwidth}
    \begin{beamercolorbox}[wd=\linewidth,sep=5pt,rounded=true]{block title}
      {\fontsize{6.8}{8}\selectfont\color{ubHead}SUB-CPMK MENUJU 1.1}\\[2pt]
      {\footnotesize Menjelaskan struktur \emph{game loop} (input,
       \emph{update}, \emph{render}) dan peran manajemen waktu dalam
       permainan.}
    \end{beamercolorbox}
    \vspace{0.16cm}

    {\fontsize{7}{8.5}\selectfont\bfseries\color{ubAccent}INDIKATOR}\par
    \vspace{0.06cm}
    \begin{itemize}
      \item Menjelaskan urutan tahapan input--\emph{update}--\emph{render}
            dengan benar.
      \item Mengidentifikasi peran \texttt{pygame.time.Clock} dan pembatasan
            FPS.
      \item Membedakan \emph{fixed timestep} dan \emph{variable timestep}
            secara konseptual.
    \end{itemize}
  \end{column}
  \begin{column}{0.45\textwidth}
    {\fontsize{7}{8.5}\selectfont\bfseries\color{ubAccent}TUJUAN BAB III}\par
    \vspace{0.08cm}
    \begin{itemize}
      \item Menyiapkan lingkungan kerja Python dan memeriksa kesiapannya
            secara mandiri.
      \item Memasang Pygame serta menjelaskan kedudukannya di antara program
            dan perangkat keras.
      \item Menyusun struktur folder proyek dan mengelola pekerjaan dengan
            Git.
      \item Menulis, menjalankan, dan mengubah program Pygame pertama.
      \item Membaca pesan kesalahan Python dan menentukan penyebabnya.
      \item Membedakan peringatan (\emph{warning}) dari kesalahan
            (\emph{error}).
    \end{itemize}
  \end{column}
\end{columns}
\end{frame}

% --------------------------------------------------------- peta
\petapertemuan{1}
  {Pertemuan pertama (orientasi dan kontrak belajar)}
  {Materi 1 \textperiodcentered\ Pengantar dan persiapan lingkungan}
  {Materi 2 \textperiodcentered\ Game Loop, Event Handling, dan Clock}
  {Lingkungan kerja siap dipakai, program Pygame pertama berjalan, dan
   repositori Git berisi penyimpanan awal.}

% ========================================================= bagian 1
\section{Dari Teori ke Praktik}
\sectionframe{01}{Dari Teori ke Praktik}
  {Tiga hasil yang harus ada, bukan sekadar dipahami}

\begin{frame}{Tiga Hasil yang Harus Diperoleh}
\awalan{Bab ini bukan bab teori. Ia selesai apabila tiga hal berikut benar-benar
ada di komputer Anda.}
\begin{itemize}
  \item \textbf{Lingkungan kerja yang siap.} Python dan Pygame terpasang,
        dapat dipanggil, dan Anda mampu memeriksanya sendiri.
        \begin{itemize}
          \item Tanpa bantuan orang lain, termasuk saat berpindah komputer.
        \end{itemize}
  \item \textbf{Program pertama yang berjalan.} Jendela terbuka, teks versi
        tampil, dan program menunggu sampai jendela ditutup.
  \item \textbf{Tempat kerja yang tertata.} Folder proyek dengan struktur
        jelas dan repositori Git sejak minggu pertama.
\end{itemize}
\vspace{0.2cm}
\begin{kotak}
Ketiganya dipakai terus sampai akhir semester. Praktikum minggu-minggu
berikutnya menumpuk berkas pada struktur folder yang sama, dan seluruh nilai
portofolio diambil dari repositori yang Anda buat sekarang.
\end{kotak}
\end{frame}

\begin{frame}{Konvensi Penulisan Kode pada Buku Ini}
\awalan{Satu kebiasaan perlu diketahui sebelum menulis kode pertama.}
\begin{itemize}
  \item \textbf{Nama variabel, fungsi, dan kelas berbahasa Inggris} ---
        \texttt{screen}, \texttt{check\_pygame}, \texttt{parse\_arguments} ---
        agar sejalan dengan dokumentasi resmi Pygame.
  \item \textbf{Komentar, keterangan fungsi, dan keluaran program berbahasa
        Indonesia} agar maksud setiap bagian kode dapat dipahami tanpa
        penerjemahan.
  \item \textbf{Istilah konsep pada teks berbahasa Indonesia}; istilah asingnya
        dicetak miring pada kemunculan pertama.
\end{itemize}
\vspace{0.2cm}
\begin{kiat}
Kebiasaan ini juga berguna di luar perkuliahan: kode yang dibaca tim atau
pewawancara kerja lazim memakai penamaan berbahasa Inggris.
\end{kiat}
\end{frame}

\begin{frame}{Batas Pembahasan Bab Ini}
\awalan{Mengetahui batas ini mencegah Anda mencari-cari bagian yang memang
belum ada.}
\begin{itemize}
  \item \textbf{Siklus permainan} --- baru dibahas pada Bab IV. Program bab ini
        menggambar sekali lalu menunggu peristiwa.
  \item \textbf{Gerak dan fisika} --- dibahas pada Bab V sampai Bab VII. Objek
        bergerak, gravitasi, dan tabrakan belum muncul di sini.
  \item \textbf{Pemasangan Unity} --- dibahas pada Bab XII, karena kebutuhan
        perangkatnya berbeda dari Python.
\end{itemize}
\vspace{0.2cm}
\begin{kotak}
Kerangka Bab I dan Bab II --- waktu, ruang, fisika, dan pengendalian objek ---
menjadi latar bab ini. Bab III memindahkan kerangka itu ke meja kerja: menyiapkan
perkakas, menyusun tempat bekerja, dan menulis program yang benar-benar berjalan.
\end{kotak}
\end{frame}

% ========================================================= bagian 2
\section{Menyiapkan Perkakas}
\sectionframe{02}{Menyiapkan Perkakas}
  {Dua perkakas wajib, dua yang sangat dianjurkan}

\begin{frame}{Empat Perkakas dan Perannya}
\begin{itemize}
  \item \textbf{Penafsir Python} (wajib) --- menjalankan kode Python. Berkas
        \texttt{.py} bukan program yang dapat dijalankan langsung oleh sistem
        operasi; ia diserahkan kepada penafsir.
  \item \textbf{Editor kode} (wajib) --- tempat menulis kode: menyoroti
        sintaksis, mengatur indentasi, menyimpan berkas \texttt{.py}.
  \item \textbf{Terminal} (sangat dianjurkan) --- menampilkan pesan kesalahan
        secara lengkap, tempat perintah pemasangan paket dan perintah Git.
  \item \textbf{Git} (wajib untuk tugas) --- mencatat riwayat perubahan berkas.
        Seluruh tugas praktikum dikumpulkan melalui repositori Git pribadi.
\end{itemize}
\vspace{0.2cm}
\begin{kotak}
Pada minggu pertama cukup Python, editor, dan Git. Unity baru diperlukan mulai
Bab XII.
\end{kotak}
\end{frame}

\begin{frame}{Tiga Syarat Editor Kode}
\awalan{Editor apa pun boleh, asalkan memenuhi tiga syarat berikut.}
\begin{enumerate}
  \item \textbf{Menyoroti sintaksis.} Warna berbeda untuk kata kunci, teks, dan
        komentar membuat kesalahan seperti tanda kutip yang belum ditutup mudah
        terlihat.
  \item \textbf{Mengatur indentasi secara otomatis.} Python sangat peka terhadap
        indentasi; editor membantu mempertahankannya konsisten.
  \item \textbf{Memudahkan penyimpanan berkas berakhiran \texttt{.py}} pada
        folder yang Anda tentukan.
\end{enumerate}
\vspace{0.2cm}
\begin{kiat}
Editor dengan terminal terpadu dan penampil pesan kesalahan menghemat waktu.
Pilihan lazim: Visual Studio Code, PyCharm Community, atau editor teks dengan
tambahan penyorot sintaksis.
\end{kiat}
\end{frame}

\begin{frame}[fragile]{Terminal: Dua Perintah yang Perlu Dikuasai}
\awalan{Hanya dua perintah yang perlu dikuasai lebih dahulu: berpindah folder
dan melihat isi folder.}
\begin{lstlisting}
cd proyek-pygame        # masuk ke folder proyek
dir                     # Windows: daftar berkas pada folder aktif
ls                      # Linux dan macOS: daftar berkas
\end{lstlisting}
\kodejelas{\texttt{cd} berarti \emph{change directory}. Kesalahan ``berkas tidak
ditemukan'' paling sering terjadi bukan karena berkasnya hilang, melainkan
karena terminal berada di folder yang keliru.}
\end{frame}

% ========================================================= bagian 3
\section{Memasang Pygame dan Memeriksa Pemasangan}
\sectionframe{03}{Memasang Pygame dan Memeriksa Pemasangan}
  {Pustaka, pengelola paket, dan skrip pemeriksaan}

\begin{frame}{Kedudukan Pygame di atas SDL}
\centering
\adjustbox{max width=0.86\linewidth,max height=0.52\textheight}{%
  % chapters/03/figures/lapisan-perkakas.tikz
% Lapisan perkakas: program Anda memanggil Pygame, Pygame memakai SDL,
% SDL berbicara dengan sistem operasi dan perangkat.
\begin{tikzpicture}[
  lapis/.style={
    rectangle, rounded corners=2pt, draw=ubPrimary, thick, fill=ubAlt,
    align=center, text width=70mm, inner sep=4pt, minimum height=10mm,
    font=\footnotesize
  },
  lapissdl/.style={lapis, draw=ubAccent, fill=white},
  lapisos/.style={lapis, draw=ubLine, fill=ubGray, text=black},
  panah/.style={-{Stealth[length=2.2mm]}, thick, draw=ubAccent},
  sisi/.style={font=\scriptsize\itshape, ubPrimary, align=left}
]

\node[lapis] (program) at (0,0)
  {\textbf{Program Anda}\\[1pt]\scriptsize kode yang Anda tulis pada folder \code{code/}};
\node[lapissdl] (pygame) at (0,-1.8)
  {\textbf{Pygame}\\[1pt]\scriptsize jendela, gambar, huruf, bunyi, dan penanganan peristiwa};
\node[lapissdl] (sdl) at (0,-3.6)
  {\textbf{SDL}\\[1pt]\scriptsize lapisan penerjemah ke perangkat keras};
\node[lapisos] (os) at (0,-5.4)
  {\textbf{Sistem operasi dan perangkat}\\[1pt]\scriptsize layar, papan tombol, tetikus, pengeras suara};

\draw[panah] (program) -- (pygame);
\draw[panah] (pygame) -- (sdl);
\draw[panah] (sdl) -- (os);

\node[sisi, right=5mm of program, text width=42mm, anchor=west]
  {perintah yang Anda tulis: \code{pygame.display.set_mode}};
\node[sisi, right=5mm of sdl, text width=42mm, anchor=west]
  {diterjemahkan ke perintah perangkat};
\end{tikzpicture}
}
\vspace{0.12cm}
\begin{kotak}
Karena berlapis, pesan kesalahan kadang menyebut SDL, bukan Pygame. Versi SDL
yang terpasang ikut menentukan kemampuan yang tersedia.
\end{kotak}
\end{frame}

\begin{frame}[fragile]{Memasang Pygame dengan pip}
\begin{lstlisting}
python -m pip install pygame

# bila pemasangan gagal karena masalah izin
python -m pip install --user pygame
\end{lstlisting}
\kodejelas{Bentuk \texttt{python -m pip} memastikan pemasangan mengikuti
penafsir yang akan menjalankan program.}
\vspace{0.15cm}
\begin{peringatan}
Perintah \texttt{pip} saja dapat menunjuk versi Python yang berbeda. Akibatnya
Pygame terpasang pada penafsir yang tidak Anda pakai, dan program tetap
melaporkan bahwa Pygame tidak ditemukan.
\end{peringatan}
\end{frame}

\begin{frame}[fragile]{Skrip Pemeriksaan Lingkungan}
\awalan{Pemasangan yang berhasil belum membuktikan lingkungan siap. Skrip ini
memeriksa penafsir yang dipakai dan kesiapan Pygame.}
\lstinputlisting[language=Python,firstline=37,lastline=44]
  {../buku_ajar/code/python/ch03_periksa_lingkungan.py}
\kodejelas{Skrip melaporkan versi Python, Pygame, dan SDL, pengandar video, serta
huruf bawaan. Kode keluar 0 berarti siap; 1 berarti ada yang perlu diperbaiki.}
\end{frame}

\begin{frame}{Masalah Pemasangan yang Sering Terjadi}
{\scriptsize
\renewcommand{\arraystretch}{1.16}
\begin{tabularx}{\textwidth}{@{}>{\raggedright\arraybackslash}p{3.4cm}
    >{\raggedright\arraybackslash}p{3.7cm}
    >{\raggedright\arraybackslash}X@{}}
\toprule
\textbf{Gejala} & \textbf{Penyebab paling sering} & \textbf{Tindakan} \\
\midrule
\texttt{python} tidak dikenali &
  Python belum terpasang atau tidak terdaftar pada PATH &
  Pasang ulang dengan ``Add python.exe to PATH'', atau pakai \texttt{py} \\
\texttt{ModuleNotFoundError} &
  Pygame belum terpasang, atau terpasang pada penafsir berbeda &
  Perhatikan baris \emph{Penafsir Python}, lalu
  \texttt{python -m pip install pygame} \\
\texttt{Access is denied} saat memasang &
  Tidak ada izin menulis pada folder instalasi Python &
  Tambahkan pilihan \texttt{--user} \\
\texttt{No available video device} &
  Dijalankan pada lingkungan tanpa layar &
  Kerjakan di laboratorium; untuk uji tanpa layar setel
  \texttt{SDL\_VIDEODRIVER} bernilai \texttt{dummy} \\
\bottomrule
\end{tabularx}}
\vfill
\begin{kotak}
Jendela yang terbuka lalu langsung tertutup biasanya dijalankan dengan pilihan
\texttt{--auto-close}: periksa argumen pada baris perintah.
\end{kotak}
\end{frame}

% ========================================================= bagian 4
\section{Struktur Folder Proyek dan Pengelolaan Versi}
\sectionframe{04}{Struktur Folder dan Pengelolaan Versi}
  {Tempat bekerja yang tertata dan riwayat pekerjaan}

\begin{frame}{Struktur Folder Proyek}
\centering
\adjustbox{max width=0.98\linewidth,max height=0.55\textheight}{%
  % chapters/03/figures/struktur-proyek.tikz
% Struktur folder proyek yang dipakai sepanjang semester.
% Pohon digambar mendatar agar nama berkas yang panjang tetap muat.
\begin{tikzpicture}[
  dir/.style={
    rectangle, rounded corners=1pt, draw=ubPrimary, thick, fill=ubAlt,
    inner sep=3.5pt, font=\ttfamily\footnotesize, text=ubPrimary
  },
  berkas/.style={
    rectangle, rounded corners=1pt, draw=ubLine, fill=white,
    inner sep=3.5pt, font=\ttfamily\footnotesize
  },
  grow'=right,
  level 1/.style={sibling distance=9mm, level distance=40mm},
  level 2/.style={sibling distance=7mm, level distance=40mm},
  edge from parent/.style={draw=ubAccent, thick}
]

\node[dir, anchor=west] (akar) at (-6.2,0) {proyek-pygame}
  child { node[dir] {code}
    child { node[berkas] {ch03\_jendela.py} }
    child { node[berkas] {ch05\_gerak.py} } }
  child { node[dir] {catatan}
    child { node[berkas] {praktikum-01.md} } }
  child { node[dir] {aset} }
  child { node[berkas] {.gitignore} }
  child { node[berkas] {README.md} };

\node[font=\scriptsize\itshape, ubPrimary, align=center, text width=112mm]
  at (0,-3.3)
  {Nama folder dan berkas ditulis dengan huruf kecil dan garis bawah, tanpa
   spasi, agar aman dipakai pada semua sistem operasi.};
\end{tikzpicture}
}
\vspace{0.10cm}
\begin{columns}[T,onlytextwidth]
  \begin{column}{0.32\textwidth}
    {\scriptsize\textbf{\texttt{code}} --- seluruh berkas program, satu berkas
     per latihan.}
  \end{column}
  \begin{column}{0.34\textwidth}
    {\scriptsize\textbf{\texttt{catatan}} --- kendala dan cara penyelesaiannya;
     bahan laporan proyek akhir.}
  \end{column}
  \begin{column}{0.30\textwidth}
    {\scriptsize\textbf{\texttt{aset}} --- gambar dan bunyi; baru terisi mulai
     Bab VII.}
  \end{column}
\end{columns}
\end{frame}

\begin{frame}[fragile]{Berkas .gitignore}
\awalan{Tidak semua berkas perlu disimpan. Salinan sementara Python dapat dibuat
ulang kapan saja dan berbeda pada setiap komputer.}
\begin{lstlisting}
__pycache__/
*.pyc
.venv/
.vscode/
\end{lstlisting}
\kodejelas{Empat baris mengabaikan salinan sementara Python, berkas terjemahan
kode, lingkungan virtual, dan pengaturan editor yang bersifat pribadi.}
\end{frame}

\begin{frame}[fragile]{Tiga Tempat Penyimpanan Git}
\begin{itemize}
  \item \textbf{Folder kerja} (\emph{working directory}) --- berkas yang sedang
        Anda sunting.
  \item \textbf{Daftar siap simpan} (\emph{staging area}) --- berkas yang sudah
        ditandai untuk ikut disimpan.
  \item \textbf{Riwayat} (\emph{repository}) --- penyimpanan yang telah tercatat
        dan tidak berubah lagi.
\end{itemize}
\vspace{0.10cm}
\begin{lstlisting}
git init                  # sekali saja: memulai repositori
git status                # memeriksa keadaan berkas
git add .                 # menandai seluruh perubahan
git commit -m "praktikum 01: jendela pertama dan struktur proyek"
git log --oneline         # melihat riwayat penyimpanan
\end{lstlisting}
\kodejelas{\texttt{git add} tidak menyimpan apa pun ke riwayat; penyimpanan
sebenarnya terjadi pada \texttt{git commit}.}
\end{frame}

% ========================================================= bagian 5
\section{Program Pygame Pertama}
\sectionframe{05}{Program Pygame Pertama}
  {Jendela, teks versi, dan perulangan peristiwa}

\begin{frame}{Kerangka Program dalam Lima Langkah}
\awalan{Kelima langkah ini akan Anda temui kembali, dengan tambahan, pada
seluruh bab berikutnya.}
\begin{enumerate}
  \item \textbf{Memulai Pygame} dengan \texttt{pygame.init()}, sehingga seluruh
        bagian pustaka yang diperlukan siap dipakai.
  \item \textbf{Membuat permukaan tampilan} dengan
        \texttt{pygame.display.set\_mode()}, yang membuka jendela dan
        mengembalikan permukaan tempat Anda menggambar.
  \item \textbf{Menggambar isi jendela}: mengisi latar dengan satu warna, lalu
        menempelkan teks atau bentuk di atasnya.
  \item \textbf{Menayangkan hasil gambar} dengan \texttt{pygame.display.flip()};
        gambar yang telah disusun masih berada di memori.
  \item \textbf{Menunggu sampai program diakhiri, lalu membersihkan} dengan
        \texttt{pygame.quit()}.
\end{enumerate}
\vspace{0.15cm}
\begin{kotak}
Pada Pygame, apa pun yang digambar --- jendela, gambar yang dimuat, atau teks ---
adalah \textbf{permukaan} (\emph{surface}). Menggambar berarti menempelkan satu
permukaan ke permukaan lain dengan \texttt{blit}.
\end{kotak}
\end{frame}

\begin{frame}[fragile]{Program Pertama: Jendela dan Teks}
\lstinputlisting[language=Python,firstline=56,lastline=71]
  {../buku_ajar/code/python/ch03_jendela.py}
\kodejelas{\texttt{set\_mode} mengembalikan permukaan \texttt{screen}; warna
dinyatakan rangkap tiga 0--255; \texttt{font.render} mengubah teks menjadi
permukaan; \texttt{blit} menempelkannya; \texttt{flip} menayangkan hasilnya.}
\end{frame}

\begin{frame}{Perulangan Peristiwa, Bukan Siklus Permainan}
\begin{itemize}
  \item \texttt{pygame.event.wait()} menghentikan program sampai satu peristiwa
        (\emph{event}) terjadi, misalnya tombol ditekan atau jendela ditutup.
  \item Bila peristiwa itu \texttt{QUIT}, variabel \texttt{running} diubah
        menjadi salah dan perulangan berakhir.
  \item Tidak ada keadaan yang diperbarui dan tidak ada gambar yang disusun
        ulang di dalam perulangan ini.
  \item Akibatnya isi jendela tidak digambar kembali walaupun tertutup sebagian
        oleh jendela lain.
\end{itemize}
\vspace{0.2cm}
\begin{peringatan}
Perulangan ini adalah \textbf{perulangan peristiwa}. Siklus permainan yang
membaca masukan, memperbarui keadaan, dan menggambar ulang terus-menerus baru
dibahas pada Bab IV.
\end{peringatan}
\end{frame}

% ========================================================= bagian 6
\section{Menjalankan, Memeriksa, dan Menelusuri Kesalahan}
\sectionframe{06}{Menjalankan dan Menelusuri Kesalahan}
  {Membaca keluaran, traceback, dan membedakan peringatan}

\begin{frame}{Alur Kerja Harian}
\centering
\adjustbox{max width=0.92\linewidth,max height=0.52\textheight}{%
  % chapters/03/figures/alur-kerja.tikz
% Alur kerja harian: tulis, jalankan, amati, simpan; kembali menulis bila
% hasilnya belum sesuai harapan.
\begin{tikzpicture}[
  langkah/.style={
    rectangle, rounded corners=2pt, draw=ubPrimary, thick, fill=ubAlt,
    align=center, text width=26mm, inner sep=3.5pt, minimum height=11mm,
    font=\footnotesize
  },
  akhir/.style={langkah, draw=ubGreen, fill=ubGreenBg},
  panah/.style={-{Stealth[length=2.4mm]}, thick, draw=ubAccent},
  ulang/.style={-{Stealth[length=2.4mm]}, thick, draw=ubWarn, dashed},
  label/.style={font=\scriptsize, ubWarn}
]

\node[langkah] (tulis) at (0,0) {Tulis atau\\perbaiki kode};
\node[langkah, right=8mm of tulis] (jalan) {Jalankan\\program};
\node[langkah, right=8mm of jalan] (amati) {Amati keluaran\\dan pesan galat};
\node[akhir, right=8mm of amati] (simpan) {Simpan dan\\catat hasilnya};

\draw[panah] (tulis) -- (jalan);
\draw[panah] (jalan) -- (amati);
\draw[panah] (amati) -- (simpan);

\draw[ulang] (amati.south) -- ++(0,-10mm) -| (tulis.south);
\node[label, anchor=north] at (4.5,-1.5)
  {bila hasilnya belum sesuai harapan};
\end{tikzpicture}
}
\vspace{0.14cm}
\begin{itemize}
  \item Program dijalankan dari terminal, tepatnya dari folder tempat berkas
        program berada.
  \item Untuk mengubah ukuran jendela tanpa menyunting kode, tambahkan pilihan
        pada perintah: \texttt{python ch03\_jendela.py --width 800 --height 600}.
\end{itemize}
\begin{kotak}
Pendekkan perulangan tulis--jalankan--amati--simpan. Perubahan kecil yang
langsung dijalankan lebih mudah dilacak daripada perubahan besar yang melibatkan
puluhan baris sekaligus.
\end{kotak}
\end{frame}

\begin{frame}[fragile]{Membaca Traceback dari Baris Terakhir}
\begin{lstlisting}[language={},numbers=none]
Traceback (most recent call last):
  File "proyek-pygame\code\galat2.py", line 5, in <module>
    screen.fill(WARNA_LATAR)
                ^^^^^^^^^^^
NameError: name 'WARNA_LATAR' is not defined
\end{lstlisting}
\kodejelas{Baris terakhir memuat jenis kesalahan; baris di atasnya memuat berkas
dan nomor baris penyebabnya. Karena itu pesan dibaca dari bawah ke atas.}
\vspace{0.12cm}
\begin{peringatan}
Kesalahan nama baru terdeteksi ketika barisnya dieksekusi. Cara paling andal
menemukannya adalah menjalankan program, bukan membacanya berulang kali.
\end{peringatan}
\end{frame}

\begin{frame}{Peringatan Bukan Kesalahan}
\begin{columns}[T,onlytextwidth]
  \begin{column}{0.48\textwidth}
    \begin{block}{Kesalahan (\emph{error})}
      \begin{itemize}
        \item Menghentikan program, atau membuat program tidak dapat dijalankan
              sama sekali.
        \item Selama masih ada pesan kesalahan, program tidak akan berjalan.
        \item Contoh: \texttt{SyntaxError}, \texttt{IndentationError},
              \texttt{NameError}, \texttt{ModuleNotFoundError}.
      \end{itemize}
    \end{block}
  \end{column}
  \begin{column}{0.48\textwidth}
    \begin{exampleblock}{Peringatan (\emph{warning})}
      \begin{itemize}
        \item Memberitahukan sesuatu yang sebaiknya diketahui.
        \item Program tetap berjalan; boleh diabaikan apabila maksudnya
              dipahami.
        \item Contoh: sambutan Pygame dan peringatan \texttt{pkg\_resources}
              yang sudah usang.
      \end{itemize}
    \end{exampleblock}
  \end{column}
\end{columns}
\vspace{0.18cm}
\begin{kotak}
Yang menandakan program berjalan baik adalah keluaran program Anda sendiri:
``Jendela 640x480 dibuka.'' dan ``Program selesai.''
\end{kotak}
\end{frame}

\begin{frame}{Kesalahan yang Sering Terjadi}
{\scriptsize
\renewcommand{\arraystretch}{1.18}
\begin{tabularx}{\textwidth}{@{}>{\raggedright\arraybackslash}p{4.4cm}
    >{\raggedright\arraybackslash}p{3.6cm}
    >{\raggedright\arraybackslash}X@{}}
\toprule
\textbf{Pesan atau gejala} & \textbf{Penyebab} & \textbf{Perbaikan} \\
\midrule
\texttt{SyntaxError: '(' was never closed} &
  Ada kurung, tanda kutip, atau kurung siku yang belum ditutup &
  Periksa baris yang ditunjuk tanda panah; hitung pasangan kurung \\
\texttt{IndentationError: unexpected indent} &
  Indentasi tidak konsisten, misalnya bercampur tab dan spasi &
  Pakai empat spasi setiap tingkat; atur editor mengubah tab menjadi spasi \\
\texttt{NameError: name 'X' is not defined} &
  Nama dipakai sebelum dibuat, atau ejaannya berbeda &
  Periksa ejaan dan besar-kecil huruf \\
\texttt{ModuleNotFoundError: No module named 'pygame'} &
  Pygame belum terpasang pada penafsir yang sedang dipakai &
  Jalankan skrip pemeriksaan, lalu \texttt{python -m pip install pygame} \\
\texttt{FileNotFoundError: [Errno 2]} &
  Folder kerja terminal bukan folder tempat berkas berada &
  \texttt{cd} ke folder yang benar, lalu periksa dengan \texttt{dir} atau
  \texttt{ls} \\
\bottomrule
\end{tabularx}}
\end{frame}

% --------------------------------------------------------- aktivitas
\begin{frame}{Aktivitas di Kelas \textperiodcentered\ Praktikum 01}
\awalan{Pengalaman belajar minggu 1 pada RPS, dikerjakan berurutan di kelas lalu
diselesaikan di rumah.}
\begin{enumerate}
  \item Periksa lingkungan: jalankan \texttt{ch03\_periksa\_lingkungan.py},
        catat versi Python, Pygame, SDL, dan nama pengandar video.
  \item Susun folder \texttt{proyek-pygame} berisi \texttt{code},
        \texttt{catatan}, dan \texttt{aset}, lalu salin skrip pemeriksaan ke
        dalamnya.
  \item Ketik ulang \texttt{ch03\_jendela.py} --- jangan salin-tempel ---
        jalankan, lalu tutup jendelanya.
  \item Jalankan dengan \texttt{--width 800 --height 600}, ubah
        \texttt{WARNA\_LATAR}, dan tambahkan satu baris nama serta NIM Anda.
  \item Buat \texttt{.gitignore}, simpan pekerjaan dengan pesan yang menerangkan
        isinya, lalu tulis catatan kendala pada folder \texttt{catatan}.
\end{enumerate}
\vspace{0.12cm}
{\footnotesize\color{ubAccent}\itshape Rincian lengkap: Praktikum 01 pada Bab III
buku ajar. Kumpulkan melalui repositori Git pribadi.}
\end{frame}

% --------------------------------------------------------- rangkuman
\begin{frame}{Rangkuman}
\begin{itemize}
  \item Python dijalankan oleh penafsir; pemasangan di Windows perlu mencentang
        ``Add python.exe to PATH''.
  \item Pygame adalah pustaka di atas SDL: jendela, gambar, huruf, bunyi, dan
        masukan.
  \item \texttt{python -m pip install pygame} memasang pada penafsir yang sama
        dengan program.
  \item Skrip pemeriksaan melaporkan versi, pengandar video, dan huruf bawaan.
  \item Folder \texttt{code}/\texttt{catatan}/\texttt{aset} dengan
        \texttt{README.md} dan \texttt{.gitignore}; Git mencatat riwayat.
  \item Menggambar dan menayangkan berbeda: perubahan baru terlihat setelah
        \texttt{pygame.display.flip()}.
\end{itemize}
\end{frame}

% --------------------------------------------------------- cek
\begin{frame}{Cek Pemahaman}
\begin{enumerate}
  \item Mengapa perintah \texttt{pip} saja berisiko, sedangkan
        \texttt{python -m pip} tidak?
  \item Apa yang terjadi bila \texttt{pygame.display.flip()} dihilangkan setelah
        penggambaran?
  \item Mengapa \texttt{NameError} tidak dapat ditemukan sebelum program
        dijalankan?
\end{enumerate}
\vspace{0.35cm}
{\footnotesize\color{ubMuted}\itshape Jawab tanpa membuka catatan, lalu
bandingkan dengan pasangan di sebelah Anda.}
\end{frame}

% --------------------------------------------------------- penutup
\penutup{Tugas Praktikum 01}
  {Pertemuan berikutnya menyusun siklus input $\rightarrow$ \emph{update}
   $\rightarrow$ \emph{render}: menangkap \emph{event} \texttt{QUIT},
   membatasi laju bingkai dengan \emph{Clock}, dan membedakan \emph{fixed}
   dengan \emph{variable timestep}.}

\end{document}
